
En
appliquant la formule 2.8, on obtient :
@@M0@@
Si on pose :
@@M1@@
cette dernière relation peut encore s'écrire :
@@M2@@
Comme le terme constant de @@M3@@ est @@M4@@, et qu'on écrit des congruences modulo torsion, on peut intégrer et écrire la congruence @@M5@@. Ceci donne les congruences annoncées.

\textbf{Corollaire 4.4.} \uline{Avec les notations et les hypothèses de} 4.3, \uline{on a pour tout} @@M6@@, \uline{et tout} @@M7@@ :
@@M8@@

Le corollaire résulte immédiatement de 4.3, en utilisant l'expression des @@M9@@ comme combinaisons linéaires à coefficients entiers des @@M10@@.

\textbf{Remarque 4.5.} La démonstration même de 4.3 laisse penser que (4.3.1) doit être vrai non seulement modulo torsion, mais aussi en tant qu'égalité. Dans [RRR], l'égalité était montrée en caractéristique @@M11@@ grâce à une caractérisation géométrique des @@M12@@. Une extension de cette méthode en toutes caractéristiques serait sans doute possible si on savait définir les opérations de puissance extérieure sur les complexes parfaits (cf. XIV).

